\section{Evolutionary Graph Theory}
\label{sec:egt}

Beyond the domain of well-mixed chemostat conditions, spatial structure is ubiquitous across natural biological populations.
Evolutionary computation, by contrast, typically defaults to well-mixed conditions; that is, any individual in the population is eligible to compete and/or sexually reproduce with any other individual during the selection process. % @AML added clarification here because well-mixed is less commonly used as an EC term
% @AML I would definitely work in references to age-layered population structure, hierarchical fair competition, super explorers, etc
Outside typical well-mixed conditions, many popular evolutionary computation techniques do impose structure on their populations, such as age-layered population structure (ALPS) [citations], hierarchical fair competition [cite], MAP-Elites [cite], and many others [citations].
However, these methods are less commonly thought of as spatially structured, but instead, structure populations according to ancestral or phenotypic traits.

While EC work has often explored the consequences of incorporating spatial structure for its own sake
\citep{tomassini2005spatially,simon2008biogeographybased}, the topic of population structure most commonly arises in the context of parallel and distributed hardware resources \citep{cantupaz1999scalability}.
To harness such hardware, independent subpopulations can be instantiated across available processors --- as if ``islands'' linked by occasional migration \citep{cohoon1987punctuated,grosso1985computer,tanese1989distributed}.
The question of how best to arrange such island subpopulations has attracted substantial attention within EC \citep{whitley1999island,cantupaz1999topologies,rucinski2010impact}.
Topics of interest have included migration policy \citep{cantupaz2001migration,cantupaz1999topologies}, population size \citep{whitley1999island,goldberg1989sizing,cantupaz2001efficient}, population diversity \citep{tanese1989distributed,whitley1999island}, and which problem classes benefit from population structure \citep{lassig2013design}.

Outside of EC, great interest also exists in how spatial structure can influence natural selection.
This question is studied directly under the banner of evolutionary graph theory (EGT) \citep{lieberman2005evolutionary}.
Under an EGT framing, population structure is typically treated as a graph topology --- with an individual occupying each vertex and edges indicating possible routes for lineage spread \citep{pattni2015evolutionary}.

A central result of evolutionary graph theory is that population structure can systematically influence evolutionary outcomes: relative to a well-mixed population, some population configurations amplify selection by making beneficial mutations more likely to fix, while others instead suppress these fixation probabilities \citep{lieberman2005evolutionary,pavlogiannis2018construction,allen2017evolutionary}.
Similar work has also examined how graph topology influences the fate of deleterious mutations \citep{sharma2022suppressors}.

Importantly, EGT has shown population structure to shape not only whether a variant fixes, but also how quickly \citep{askari2015analytical}.
Surprisingly, in such scenarios, faster fixation time often comes at the cost of lower fixation certainty \citep{frean2013effect}.
Thematic parallels exist between such EGT work and EC's interest in extending duration required for mutational sweeps (``takeover time''), as a mechanism to slow search convergence \citep{goldberg1991comparative,rudolph2001takeover,giacobini2005selection,sudholt2015parallel}.

As well-known in EC \citep{wolpert1997no}, no single search configuration can be expected optimal across all possible fitness landscapes.
As such, recent work in EGT has developed techniques to engineer bespoke population structures tuned to individual optimization problems \citep{kuo2024theory,kuo2024evolutionary}.
This framework revolves around manipulation of two evolutionary properties: selection strength (``amplification factor'') \citep{kuo2024theory} and effective population size/fixation time (``acceleration factor'') \citep{kuo2024evolutionary}.
Crucially, these properties may be tuned independently --- a capability achieved by use of higher-order triangle motifs that alter fixation time without affecting fixation probability \citep{kuo2024evolutionary}.
Such principled construction refines earlier reliance on black box optimization (notably, including EC approaches) \citep{moller2019exploring}.

To test population structures constructed using their method, \citet{kuo2024evolutionary} evaluate performance on EC-style optimization tasks.
On a task in which agents evolve spatial navigation behaviors, accelerated fixation speeds solution discovery.
By contrast, to prevent premature fixation on a rugged benchmark function, ``decelerator'' motifs enhance solution quality.
\citet{kuo2024evolutionarymutation} extends analysis to consider valley-crossing behavior.
Here, amplification and acceleration factors are shown to jointly influence the probability of successfully crossing a fitness valley in a two-site landscape.

As with other facets of bioscience, potential exists to extend formalisms developed within EC by relating them to EGT.
However, given this recent intervention-oriented framing, EGT increasingly provides off-the-shelf capabilities directly relevant to EC practitioners.
EC systems fit so well within an EGT framework, in fact, that in many regards they provide an ideal model system to conduct EGT studies.
Indeed, such work is needed to characterize the limitations of classical EGT assumptions around update rules \citep{hindersin2015most}, mutation rate \citep{sharma2023selfloops}, and single-occupancy population structure \citep{broom2012general}.
Promising early work on this front has explored the influence of EGT-style population structure on longer-term evolutionary outcomes, such as life history trade-offs in digital organisms \citep{gordon2025environmental}.
Moreover, \citet{gordon2025environmental} show that populations evolving on different graph structures exhibit different rates of long-term optimization for different fitness landscapes; specifically, more constrained graph structures better facilitate crossing fitness valleys and general search space exploration, while less constrained graph structures better facilitate rapid exploitation.